% Per-sample statistics paragraph, rewritten. All numbers verified against the audit
% store: ensemble k>=16 share 5.86%; per-triplet total energy 2.8e9..6.2e9; total-energy
% max deviations base 7.0% / ft 8.0% (base+dial 8.8% / ft+dial 9.9% if wanted);
% k>=16 per-sample: base 0% in-band, slope 0.98, mean dose 0.67x; ft 88%, slope 0.78,
% mean 0.93x; ft+dial 70%, slope 0.39; base+dial 50%, slope 0.40.
% Figure caption should read k>=16, not k>=32.

\paragraph{In-distribution per sample statistics} We have shown that sampling guidance
and DDPO finetuning address part of the energy deficit of the baseline model over the
mid to high-$k$ range, and that both enable the recovery of SR samples which are more
energetically and physically consistent with the flow's ensemble statistics in the base
model's home regime. We now determine whether these improvements over the mean
statistics of the SR generations translate to the individual sample.\newline

We first justify the need for such an analysis. Figure
\ref{fig:re-1000-k16-energy-breakdown} shows that the total energy carried by individual
triplets spans the $E\approx[3,6]\times10^{9}$ range, with only $5.86\%$ of that energy
contained at wavenumbers $k\geq16$. This study places a specific emphasis on restoring
the correct mid to high-$k$ energy content of each individual sample, a quantity which
constitutes only a small part of the total sample energy, without allowing the injection
process to revert to the group mean dose. Over the full $k$ range, the base and
finetuned models both recover the total energy of each triplet accurately, with maximum
deviations $\lvert E_{x_0} - E_{GT}\rvert / E_{GT}$ of $7.0\%$ and $8.0\%$ respectively.
Restricting the range to $k\geq16$, where the conditional dose must be read from the
input rather than inherited from the large scales, both models degrade. The base model
tracks the per-sample energy variation almost perfectly, with a slope of $0.98$, but
applies a near-constant $0.67\times$ dose, so that not a single sample falls within
$20\%$ of its own truth. Its failure is purely multiplicative. The finetuned model
raises the mean dose to $0.93\times$ and places $88\%$ of its samples within the
$20\%$ tolerance, at the cost of flattening the slope to $0.78$. Samples carrying less
$k\geq16$ energy than the ensemble mean are relatively over-energised and samples
carrying more are under-energised, the signature of a reward formulation anchored to
ensemble statistics pulling every generation toward the group level.\newline
This flattening is considerably stronger when sampling guidance is applied at inference
time. The guidance gradient reaches the sample through a reward anchored to the
ensemble statistics of the regime, and it is therefore blind to the conditional dose
each sample requires, pushing every generation toward the ensemble mean. Under
guidance, the tolerance accuracy of the finetuned model falls from $88\%$ to $70\%$ and
its slope collapses from $0.78$ to $0.39$, while the base model rises from $0\%$ to
$50\%$ with a slope of $0.40$. The guided base model reaches half of its samples not by
reading their content but by dosing all of them toward the ensemble mean, which its
flattened slope makes explicit.\newline
